%! Equivalent Representations of Polynomial and Rational Expressions — Unit 1, Lesson 1.11 — Warm-Up
\documentclass[10pt]{article}
\usepackage{precalculus-article}
\usepackage{precalculus-boxes}
\newcommand{\TallMath}[1]{$\displaystyle #1\rule[-1.4em]{0pt}{3.2em}$}

\begin{document}

\pageheader{Unit 1, Lesson 1.11}{Warm-Up}
\namedateperiod

\vspace{0.12in}

\begin{enumerate}[itemsep=14pt]

\item Factor the polynomial $p(x) = x^3 - x^2 - 6x$ completely.

      \vspace{0.08in}
      $p(x) = $ \blank{7cm}

      Zeros of $p$: \blank{5cm}

\item A rational function is given: $\displaystyle r(x) = \dfrac{x^2 - 9}{(x-3)(x+5)}$.

      \begin{enumerate}[label=(\alph*), itemsep=6pt]
        \item Does $r$ have a hole, a vertical asymptote, or both at $x = 3$?
              Circle: \quad \textbf{Hole} \quad / \quad \textbf{Vertical Asymptote} \quad / \quad \textbf{Both}

        \item Identify any vertical asymptote(s). \blank{4cm}
      \end{enumerate}

\item The rational function $\displaystyle h(x) = \frac{3x^4 - x + 7}{2x^2 + 1}$ is written in
      standard form. Describe the end behavior of $h(x)$ as $x \to +\infty$.

      \vspace{0.06in}
      End behavior: as $x \to +\infty$, $h(x) \to$ \blank{3cm}

\end{enumerate}

\end{document}
